% Chapter 2 exercises. Source PDF pages 80--81 / printed pages 66--67.
\subsection*{习题}
\label{sec:ch02-exercises}
\setcounter{exerciseitem}{0}

\begin{MBExercise}
\label{ex:ch02-exercise-01}
证明例~\ref{ex:ch02-hilbert-subspaces} 的 (ii).
\end{MBExercise}

\begin{MBExercise}
\label{ex:ch02-exercise-02}
若 $M$ 是子空间, 证明 $(M^\perp)^\perp=M$.
\end{MBExercise}

\begin{MBExercise}
\label{ex:ch02-exercise-03}
证明命题~\ref{prop:ch02-orthogonal-complement-properties} 的 (1) 和 (3).
\end{MBExercise}

\begin{MBExercise}
\label{ex:ch02-exercise-04}
证明平行四边形法则, 即定理~\ref{thm:ch02-parallelogram-law}.
\end{MBExercise}

\begin{MBExercise}
\label{ex:ch02-exercise-05}
设 $P_M^\perp$ 是 \eqref{eq:ch02-pm-perp-definition} 中定义的算子. 证明 $P_M^\perp=P_{M^\perp}$.
\end{MBExercise}

\begin{MBExercise}
\label{ex:ch02-exercise-06}
证明注~\ref{rem:ch02-ritz-method} 中的断言.
\end{MBExercise}

\begin{MBExercise}
\label{ex:ch02-exercise-07}
证明压缩映射总是连续的(参见引理~\ref{lem:ch02-contraction-mapping}).
\end{MBExercise}

\begin{MBExercise}
\label{ex:ch02-exercise-08}
证明 Lax--Milgram 定理~\ref{thm:ch02-lax-milgram} 的证明中的映射 $u\mapsto Au$ 是线性映射 $V\longrightarrow V'$.
\end{MBExercise}

\MathBookSourcePage{81}
\begin{MBExercise}
\label{ex:ch02-exercise-09}
证明 Lax--Milgram 定理保证的解 $u$ 满足
\[
\MathBookFormulaID{formula-ch02-exercise-09-stability}
\|u\|_V\leq\frac1\alpha\|F\|_{V'}
\]
(参见注~\ref{rem:ch02-lax-milgram-stability}).
\end{MBExercise}

\begin{MBExercise}
\label{ex:ch02-exercise-10}
对微分方程 $-u''+ku'+u=f$, 求一个 $k$ 值, 使某个 $v\in H^1(0,1)$ 满足 $a(v,v)=0$ 但 $v\not\equiv0$(参见 \eqref{eq:ch02-large-k-differential-equation}).
\end{MBExercise}

\begin{MBExercise}
\label{ex:ch02-exercise-11}
设 $a(\cdot,\cdot)$ 是 Hilbert 空间 $V$ 的内积. 对 $F\in V'$ 和 $V$ 的任意(闭)子空间 $U$, 证明下述两个命题等价:
\begin{enumerate}
\item[(a)] $u\in U$ 满足 $a(u,v)=F(v)$, $\forall v\in U$;
\item[(b)] $u$ 在 $v\in U$ 上极小化 $\frac12a(v,v)-F(v)$.
\end{enumerate}
即证明一个命题中的存在性蕴含另一个命题中的存在性. (提示: 展开 $a(u+\epsilon v,u+\epsilon v)=\cdots$.)
\end{MBExercise}

\begin{MBExercise}
\label{ex:ch02-exercise-12}
设
\[
\MathBookFormulaID{formula-ch02-exercise-12-form-space}
a(u,v)=\int_0^1(u'v'+u'v+uv)\,dx,
\qquad
V=\{v\in W_2^1(0,1):v(0)=v(1)=0\}.
\]
证明对所有 $v\in V$,
\[
\MathBookFormulaID{formula-ch02-exercise-12-identity}
a(v,v)=\int_0^1\bigl((v')^2+v^2\bigr)\,dx.
\]
(提示: 写成 $vv'=\frac12(v^2)'$.)
\end{MBExercise}

\begin{MBExercise}
\label{ex:ch02-exercise-13}
设 $a(\cdot,\cdot)$, $V$, $U$ 如习题~\ref{ex:ch02-exercise-11}. 对 $g\in V$, 定义 $U_g=\{v+g:v\in U\}$(注意 $U_0=U$). 对任意 $g\in V$, 证明下述命题等价:
\begin{enumerate}
\item[(a)] $u\in U_g$ 满足 $a(u,v)=0$, $\forall v\in U_0$;
\item[(b)] $u$ 在所有 $v\in U_g$ 上极小化 $a(v,v)$.
\end{enumerate}
\end{MBExercise}

\begin{MBExercise}
\label{ex:ch02-exercise-14}
证明 $\mathcal D(\Omega)$ 在 $L^2(\Omega)$ 中稠密. (提示: 使用习题~\ref{ex:ch01-exercise-39} 证明 $\mathcal D(\Omega)^\perp=\{0\}$.)
\end{MBExercise}

\begin{MBExercise}
\label{ex:ch02-exercise-15}
设 $H$ 是 Hilbert 空间, $v\in H$ 任意. 证明
\[
\MathBookFormulaID{formula-ch02-exercise-15-dual-characterization}
\|v\|_H=\sup_{0\neq w\in H}\frac{(v,w)_H}{\|w\|_H}.
\]
(提示: 应用 Schwarz 不等式 \eqref{eq:ch02-schwarz-inequality}, 并考虑 $w=v$.)
\end{MBExercise}
